\scanpage{11}

\uline{Prop. \unsure{1.11|1.10}.} Soient $\mathcal E,\mathcal E'$ deux topos multigaloisiens, $\Pi_1(\mathcal E),\Pi_1(\mathcal E')$ leurs groupoïdes fondamentaux. Le foncteur évident
\[
\underline{\mathrm{Hom}}_{\mathrm{top}}(\mathcal E,\mathcal E')
\longrightarrow
\underline{\mathrm{Hom}}\bigl(
\underset{C}{\underline{\Pi_1(\mathcal E)}},
\underset{C'}{\underline{\Pi_1(\mathcal E')}}\bigr)
\]
est une équivalence des catégories, on trouve une équivalence quasi-inverse en composant
\[
\underline{\mathrm{Hom}}(\deleted{\illegible{1}}C,C')
\longrightarrow
\begin{array}[t]{c}
\underline{\mathrm{Hom}}_{\mathrm{top}}(\widehat C,\widehat {C'})\\
\cap\\
\underline{\mathrm{Hom}}(\widehat {C'},\widehat C)
\end{array}
\simeq
\underline{\mathrm{Hom}}_{\mathrm{top}}(\mathcal E,\mathcal E').
\]
\inserted{(\unsure{ou|car} $\widehat C\simeq\mathcal E$, $\widehat {C'}\simeq\mathcal E'$)}

\deleted{\illegible{1}}
\uline{1.12} Explicitation du cas où $\mathcal E\ \mathcal E'$ sont $0$-connexes et ponctués, donc donnés comme $\mathcal E\simeq\mathrm{Ens}(G)$, $\mathcal E'\simeq\mathrm{Ens}(G')$

\dots

\margin{10}

\scanpage{12}
\mspage{4}

2. \unsure{Application|Applications} aux revêtements des topos.

\uline{Thm 2.1.} Soit $\mathcal E$ un topos loc. connexe (i.e. dont tout objet est somme d'objets $0$-connexes) et loc. simplement connexe (i.e. admettant un syst. de générateurs \unsure{qui} sont $1$-connexes).
\margin{peut-être faut-il \unsure{supposer|supposé} que $\mathcal E$ ait ``assez de points'' i.e. assez de foncteurs fibres\dots}
\deleted{\illegible{6}}\inserted{Alors la catégorie $\mathcal E_{lc}$} des objets loc\textsuperscript{t} constants de $\mathcal E$ est \inserted{un topos} multigaloisien\deleted{\unsure{ne}}, et l'inclusion
\[
(2.1.1.)\qquad \mathcal E_{lc}\hookrightarrow\mathcal E
\]
commute aux $\varprojlim$ finies (NB en fait, sans hyp. sur le topos $\mathcal E$, $\mathcal E_{lc}$ est stable par $\varprojlim$ finies) \uuline{et} aux $\varinjlim$ quelc.

\uline{Déf. 2.2} On dénote ce topos par $\mathcal E_{[1]}$, on l'appelle l'enveloppe multigaloisienne de $\mathcal E$, et le morphisme de topos transposé de l'inclusion (2.1.1) :

\margin{11}

\scanpage{13}

\[
\mathcal E\longrightarrow\mathcal E_{[1]}
\]
prend le nom de \uline{morphisme canonique}.

\uline{NB} C'est l'équivalent en théorie des topos de l'opération de ``tuage des $\Pi_i$ pour $i\geq 2$''. On peut définir aussi \inserted{un} \deleted{\unsure{$\mathcal E_\infty$}} $\mathcal E_{[0]}$ et une suite
\[
\mathcal E\unsure{\longrightarrow|-}\mathcal E_{[1]}\unsure{\longrightarrow|-}\mathcal E_{[0]}
\]
[$\mathcal E_0$ est le topos \uline{discret} défini par $\Pi_0(\mathcal E)$, qui a un sens satisfaisant dès que $\mathcal E$ loc. connexe\dots]. Moyennant des hyp. convenables sur $\mathcal E$ (du type ``locale contractibilité''), on doit pouvoir définir les $\mathcal E_{[i]}$ pour $\forall\ i\in\mathbb N$, et des morphismes canoniques
\[
\mathcal E\longrightarrow\cdots\quad
\mathcal E_{[i]}\unsure{\longrightarrow|-}\mathcal E_{[i-1]}\longrightarrow\cdots\quad
\mathcal E_{[1]}\longrightarrow\mathcal E_0.
\]

\margin{12}

\scanpage{14}
\mspage{5}

2.3. Le \uline{groupoïde fondamental} de $\mathcal E$ se définit comme ayant pour objets les points de $\mathcal E$ (qui induisent des points de $\mathcal E_{[1]}$ grâce à $\mathcal E_{[1]}\longrightarrow\mathcal E$), et comme morphismes les \deleted{\unsure{points}} morphismes \uline{de points de} $\mathcal E_{[1]}$. On a donc des foncteurs canoniques
\[
\underline{\mathrm{Pt}}(\mathcal E)
\xrightarrow{\alpha}\Pi_1(\mathcal E)
\overset{\beta}{\overset{\approx}{\longrightarrow}}
\mathrm{Pt}(\mathcal E_{[1]})
\overset{\text{déf}}{=}\Pi_1(\mathcal E_{[1]})
\]
où $\beta$ est une équivalence (mais pas surjectif sur les objets), et où bien sûr $\alpha$ n'est pas nécessairement une équivalence, ni $\overline{\mathrm m}$ pl. fidèle, ou seulement fidèle. P.ex. si $\mathcal E$ est $1$-connexe i.e. $\Pi_1(\mathcal E)$ équivalent à la cat. ponctuelle, il ne s'ensuit pas nécessairement que les Hom \unsure{dans|de} $\mathrm{Pt}(\mathcal E)$ soient tous de cardinal $\leq 1$ !

\margin{13}

\scanpage{15}

Comme un p\textsuperscript{t} $P$ de $\mathcal E$ définit un p\textsuperscript{t}. (noté encore $P$ par abus) de $\mathcal E_{[1]}$, on peut donc définir le \uline{rev. universel de $\mathcal E$} basé en ce point, comme un objet \inserted{$S$} $1$-connexe de $\mathcal E_{[1]}$ -- il est caractérisé \inserted{dans $\mathcal E$} par le fait d'être loc. constant, $1$-connexe, et muni d'un relèvement
\[
P\longrightarrow\mathcal E_{/S}.
\]

Mais comme $\alpha$ n'est pas une équivalence de cat. (bien qu'il soit \deleted{\illegible{1}} \unsure{ess.} surj. si \unsure{on} suppose que $\mathcal E$ a suffisamment de pts\dots) on ne peut \uline{pas} dire que $\forall$ rev. univ. de $\mathcal E$ soit défini par un p\textsuperscript{t} de $\mathcal E$, défini à isom. unique près\dots

\margin{14}

\scanpage{16}
\mspage{6}

3. Variantes ``pro-multigaloisiennes'', resp. profinies (en se bornant, pour simplifier, au cas des topos loc. connexes\dots).

4. Complément-remords sur les cat. multigaloisiennes, qui précise l'intention que pour un topos $\mathcal E$, la donnée d'un objet $S\in\mathcal E$ définit un topos induit $\mathcal E_{/S}\longrightarrow\mathcal E$, et que $S$ se reconstitue à isom. unique près par la connaissance du topos induit en tant que topos \uline{au dessus de} $\mathcal E$. Ici, \deleted{\illegible{1} $\mathcal E_{/S}$ \unsure{est}} \inserted{$\mathcal E$ étant} multigaloisien, $\mathcal E_{/S}$ aussi -- et il se pose la question quand un morphisme de topos \inserted{\unsure{multigaloisiens|multigaloisien}} $\mathcal E'\longrightarrow\mathcal E$ peut être considéré comme

\margin{15}

\scanpage{17}

un morphisme d'induction. Or si $C=\Pi_1(\mathcal E)$, $C'=\Pi_1(\mathcal E')$, la donnée de $\mathcal E'\unsure{\longrightarrow|-}\mathcal E\deleted{\unsure{'}}$ équivaut à la donnée d'un foncteur $C'\longrightarrow C\deleted{'}$. On trouve que $\mathcal E'\unsure{\longrightarrow|-}\mathcal E\deleted{\unsure{'}}$ est un morphisme d'induction \unsure{ssi} $C'\longrightarrow C\deleted{'}$ est \uline{fidèle}. Ainsi, on trouve une équivalence entre la catégorie \inserted{$\mathcal E$} (des objets $S$ de la cat. multigaloisienne $\mathcal E\simeq\widehat C$, \deleted{\illegible{1}} où $C$ est un groupoïde, \uline{quelconque si on y tient})
\margin{un peu vif !}
et la cat. dont les objets sont les ``groupoïdes $C'$ au dessus de $C$'', avec un foncteur structural $C'\longrightarrow C$ \uline{fidèle}, les morphismes de $C'_1$ dans $C'_2$ étant les \inserted{\uuline{classes d'isomorphie \unsure{d'un|d'une}}} \deleted{\illegible{1}} couples $(f,\alpha)$ d'un foncteur $f:C'_1\unsure{\longrightarrow|-}C'_2$ et d'un

\margin{16}

\scanpage{18}
\mspage{7}

isomorphisme de foncteurs $\alpha:p_2f\simeq p_1$
\diagram{%
\begin{tikzcd}
C'_1 \arrow[rr,no head,"f"] \arrow[dr,"p_1"',""{name=pone}]
  & & C'_2 \arrow[dl,"p_2",""{name=ptwo}] \\
  & C &
\arrow[from=pone,to=ptwo,bend left=25,"\overset{\alpha}{\sim}"]
\end{tikzcd}}
\margin{Préciser les isomorphismes entre couples $(f,\alpha)$ et \deleted{\unsure{couples}} $(g,\beta)$\dots}

Dans le cas où p.ex. $C$ est la catégorie réduite à un seul objet, un groupe d'autom. $G$, cette description de la catégorie $\mathcal E=\mathrm{Ens}(G)$ est évidemment un peu lourde, mais elle s'insère bien dans certains contextes plus bas.

Ainsi, si $k$ est un corps de base, le \deleted{\unsure{$\mathcal E$}} catégorie $\mathcal E$ des schémas étales sur $k$ est décrit, en termes d'une clôture sépar. \unsure{$k$|le} $k_s$ de $k$ et du groupe profini $\Gamma=\mathrm{Gal}(k_s/k)$, comme les groupoïdes \inserted{profinis} au dessus du groupoïde profini $(\mathrm{pt},\Gamma)$\dots Nous voulons insérer cette description dans une ``description'' ``galoisienne'' \unsure{des|de} \inserted{[certains]} schémas [lisses \inserted{quasi-}projectifs de dim $\leq 1$] sur $k$, du moins si $k$ corps de type fini sur $\mathbb Q$.

\margin{17}

\scanpage{19}

5. Soit $K$ une extension, de type fini de $\mathbb Q$, et choisissons une clôture algébrique $\overline K$ de $K$.
\margin{On pose $\Gamma=\mathrm{Gal}(\overline K/K)$}
\inserted{5.1} Nous considérons des couples $(X,S)$, où : \deleted{\unsure{$X$}}

a) $X$ est un schéma projectif et lisse sur $K$, de dim $\leq 1$ ;

b) $S$ est un sous-schéma fini \inserted{réduit} de $X$, \inserted{(alors \illegible{1} \unsure{étale} sur $K$)}, contenu \deleted{\illegible{2}} \inserted{(dans)} \inserted{la réunion} \unsure{des} composantes irréductibles de dim $1$ de $X$.
\margin{(donc fini étale sur $\unsure{k|K}$)}

Les morphismes $(X',S')\unsure{\longrightarrow|-}(X,S)$ seront par définition les morphismes de schémas
\[
f:X'\longrightarrow X
\]
tels que
\[
S'=f^{-1}(S)_{\mathrm{red}}
\]
i.e. tels que $\mathrm{supp}\,S'=f^{-1}(\mathrm{supp}\,S)$.

Nous cherchons une ``description galoisienne'' de cette catégorie, ou tout au moins d'une \inserted{sous-}catégorie pleine \inserted{$\mathcal V_K$} que nous allons

\margin{18}

\scanpage{20}
\mspage{8}

définir maintenant.

\uline{Lemme} 5.2. Soit $\Omega$ un corps alg. clos, $X$ une \deleted{schéma} courbe alg. proj. \inserted{lisse}, connexe sur $\Omega$, $S$ une partie finie de $X(\Omega)$, \inserted{$g$ le genre de $X$, $n=\mathrm{card}\,S$.} $\deleted{\unsure{X_0}}\inserted{U}=X\setminus S$, Conditions équivalentes

a) $\Pi_1(U)$ non abélien

b) $\mathrm{Aut}(U)$ fini

c) Pour $\forall$ schéma \inserted{connexe \deleted{de t.f.}} réduit $X$ \inserted{de t.f.} sur $\Omega$, l'ens des morphismes \uline{non constants} de $X$ dans $U$ est fini

d) On est \deleted{\unsure{$X$}} dans l'un des trois cas suivants : $1^\circ)$ \deleted{genre} $g\geq 2$ \quad $2^\circ)$ $g=1$, $n\geq 1$ \quad $3^\circ)$ $g=0$, $n\geq 3$.

\margin{e) (si $\Omega\unsure{\in|\subset}\mathbb C$) Le rev. universel de $X(\mathbb C)\setminus S(\mathbb C)$ est isomorphe au demi plan \uline{de} Poincaré.}
\margin{\unsure{?!|??} f) (si $\Omega=\overline{\mathbb Q}$ \unsure{$S\neq\emptyset$}) Le rev. universel de $X\setminus S=U$ est isomorphe à celui de $\mathbb P^1_\Omega\setminus\{0,\unsure{1},\infty\}$.}

\uline{Déf. 5.3.} On dit alors que $(X,S)$ est anabélique. \deleted{Si $X$ est un schéma}

Comme cette condition est (par d) p.ex) invariante par extension du corps de base alg. clos, on étend cette définition au cas d'un couple $(X,S)$, avec $(X,S)$ \deleted{projectif lisse de dim} \inserted{comme dans (5.1)} (\uline{NB} on regarde séparément les comp. conn. de $X_{\overline K}$\dots).

\margin{19}
